\documentclass[11pt, a4paper, oneside]{article}
\usepackage{amsmath, amsthm, amssymb, bm, color, framed, graphicx, hyperref, mathrsfs}
\usepackage{geometry}
\usepackage{fancyhdr}
\usepackage{enumitem}
\usepackage{lastpage}
\usepackage{booktabs}
\usepackage[raggedright]{titlesec}
\titleformat{\section}
  {\normalfont\bfseries}
  {\thesection}
  {1em}
  {}

\geometry{a4paper, left=2.5cm, right=2.5cm, top=2.5cm, bottom=2.5cm}

\setlength{\parindent}{0pt}

\fancypagestyle{plain}{
    \fancyhf{}
    \fancyhead[L]{Assignment 3}
    \fancyhead[R]{CS224n Natural Language Processing}
    \fancyfoot[C]{\thepage}
}

\pagestyle{plain}

\title{\textbf{Assignment 3}}
\author{}
\date{\empty}

\begin{document}

\maketitle
\section{Neural Machine Translation with RNNs (45 points)}
\begin{enumerate}[label=(g), left=0pt, labelsep=1em, itemsep=1em]
    \item \textit{enc\_mask} marks padding tokens with 1s and other non-padding tokens 
    with 0s. And the attention scores of tokens where \textit{enc\_mask} is 1 are set 
    to $-\infty$. Therefore, with \textit{enc\_mask}, model does not attend to padding tokens 
    since their attention values are 0 after softmax.

    Such masks are necessary for NMT models to attend to meaningful tokens of a sentence, rather
    than padding tokens that are manually added to support parallel processing.
\end{enumerate}
\end{document}